% NOETHER PAPER 4 — KOREAN TRANSLATION-PRODUCER DRAFT — UNCHECKED
% Unit: T07-U36; current-authority whole lines 4241--4263
% Source slice: 2,543 LF UTF-8 bytes; SHA-256 FDAEB00ABCC919EB60EB1EDAF7FA1A7D3F4C34A5DC652A5BD5FCCD211884561B
% Pointer: NOETH-DE-AUTH-v004-20260804; SHA-256 A1C62FDACAA34DFC1B806DC18258F2E732539F7AE9D85AA4BD9E1067B8749D9F
% This is producer translation only: no source/Korean/formula review, compilation, or rendering.

정규형의 계가 형식열 $f$와 동치임을 보이기 위해 이제 남은 것은, 형식열의 모든 형식을 1.에서 정한 의미에서 정규형들의 편극들에 따라 전개할 수 있음을 보이는 것뿐이다. $\Pi$를 다음 편극 연산들의 결합식이라 하자.
\srcnumdisplay[4.2em]{(53.)}{%
 \left(\frac{\partial}{\partial\xi^{(i)}}\middle|\xi^{(k)}\right),
 \qquad (i\ne k;\ i=1,2,\ldots,\rho;\ k=1,2,\ldots,\rho)}
또한 $\Pi^\sigma$를 $i=\sigma$이고 $k<\sigma$인 특수 연산 (53.)들의 결합식이라 하자. 열 $\xi^{(\rho)}$에 관해 차수가 $k_\rho$인, $\rho$개의 열 $\xi^{(1)},\ldots,\xi^{(\rho)}$의 형식 $F$에 대해서는 다음 전개가 성립한다\srcfn{*)}{F. Mertens, \emph{\"Uber eine Formel der Determinantentheorie}. Sitzungsber. d. Ak. d. Wiss., Wien, math.-naturw. Kl., Bd. 91. 2. Abt. (1885). Formel 20).}:
\srcnumdisplay[4.2em]{(54.)}{%
 F=\sum_{\alpha=0}^{k_\rho}\Pi F_\alpha,
 \qquad(\rho\leqq n)}
여기서 $F_\alpha$는 마지막 열 $\xi^{(\rho)}$에 관해 $\alpha$차이고, $F$에 다음 불변 미분 연산
\srcnumdisplay[4.2em]{(55.)}{%
 \left(\frac{\partial}{\partial\xi^{(1)}}\cdots\frac{\partial}{\partial\xi^{(\rho)}}\middle|\xi^{(1)}\cdots\xi^{(\rho)}\right)^\alpha}
및 $\Pi^\rho$를 적용하여 생긴다. $F_\alpha$를 생성할 때 과정 (55.)를
\srcnumdisplay[4.2em]{(56.)}{%
 \left(\frac{\partial}{\partial\xi^{(1)}}\cdots\frac{\partial}{\partial\xi^{(\rho)}}\middle|p_\rho\right)^\alpha,}
로 바꾸면, $\rho-1$개의 열 $\xi^{(1)},\ldots,\xi^{(\rho-1)}$만을 가진 형식 $F_\alpha(p_\rho)$가 생기며, 여기에 다시 (54.)를 적용할 수 있다. 이 절차를 계속하면 형식 $G(p_\rho,p_{\rho-1},\ldots,p_1)$에 이른다. 이 형식들로부터 (54.)에 따른 $F$의 전개를 얻으려면, 각 열 $p_\sigma$를 $\xi^{(1)},\ldots,\xi^{(\sigma)}$로 바꾸고 생기는 형식들에 편극 과정 $\Pi$ (53.)를 적용해야 한다. 그러면 (48.)과 비교할 때 변환된 계수 $(G)$들의 합을 얻는다. $\rho=n$이면 첫 단계에서 과정 (55.)를 그대로 둔다. $c$가 수치 상수들을 뜻한다고 하고, 더 나아가 다음과 같이 줄여 쓰자.
\srcnumdisplay[4.2em]{(57.)}{%
 (\xi^{(1)}\xi^{(2)}\cdots\xi^{(n)})=\Delta;
 \qquad
 \left(\frac{\partial}{\partial\xi^{(1)}}\frac{\partial}{\partial\xi^{(2)}}\cdots\frac{\partial}{\partial\xi^{(n)}}\right)=\nabla,}
그러면 $\rho=n$인 경우에는
\srcnumdisplay[4.2em]{(58.)}{%
 F=\sum c\Delta^\alpha(G),}
을 얻는 반면, $\rho<n$이면 오른쪽에는 $\alpha=0$인 항, 곧 $\sum c(G)$만 나타난다.
